% Generated draft: requires mathematical and source review.
\documentclass[11pt]{article}
\usepackage[margin=25mm]{geometry}
\usepackage{fontspec}
\setmainfont{Latin Modern Roman}
\usepackage{amsmath,amssymb,mathtools,unicode-math}
\setmathfont{Latin Modern Math}
\usepackage{enumitem,needspace,xurl,hyperref}
\setlength{\parindent}{0pt}
\setlength{\parskip}{5pt}
\setlength{\emergencystretch}{3em}
\newcommand{\N}{\mathbb{N}}
\newcommand{\Z}{\mathbb{Z}}
\newcommand{\Q}{\mathbb{Q}}
\newcommand{\R}{\mathbb{R}}
\newcommand{\C}{\mathbb{C}}
\newcommand{\F}{\mathbb{F}}
\newcommand{\A}{\mathbb{A}}
\newcommand{\PP}{\mathbb{P}}
\newcommand{\E}{\mathbb{E}}
\newcommand{\eps}{\varepsilon}
\newcommand{\dd}{\,\mathrm{d}}
\begin{document}
\section*{UnsolvedMath ID: 30000276}
\textbf{Source problem code:} catalogue entry
\subsection{Definitions and mathematical statement}
Let \(E/\Q\) be the particular elliptic curve intended in the source, and for \(m\in\Z\) put
\[
K_m=\Q(\sqrt[3]{m}).
\]
A literal meaning of acquiring new points is strict enlargement under base change:
\[
E(K_m)\setminus E(\Q)\ne\varnothing.
\]
The classification problem is to determine
\[
\mathcal N_E=\{m\in\Z:E(K_m)\ne E(\Q)\}.
\]
Here rational points mean points rational over \(K_m\), and the inclusion of \(E(\Q)\) is the natural one. The curve's equation, conductor, or any other identifying data are not supplied, so \(\mathcal N_E\) cannot be specialized to a definite curve.

The title and assessment instead emphasize rank growth. That would ask for the potentially different set
\[
\mathcal R_E=\{m\in\Z:\operatorname{rank}E(K_m)>\operatorname{rank}E(\Q)\},
\qquad
\operatorname{rank}E(K)=\dim_{\Q}(E(K)\otimes_{\Z}\Q).
\]
Rank growth implies new points, but strict point-group enlargement could concern torsion or finite-index enlargement rather than rank. Thus the two targets must not be identified without a source-specific argument. No claim here asserts that either phenomenon occurs for the unspecified curve.

If \(m\) is an integer cube, including zero, then \(K_m=\Q\); for noncubes the extension has degree three. Different integers can define the same field, so an integer classification and a classification of distinct field extensions have different indexing. The source imposes no cube-free, positivity, or nonzero restriction on \(m\), and none is added.

The supplied assessment checked 2026-08-21 says that the passage is too truncated to identify a solved finite instance and that the corresponding general rank-growth question remains open. The corrected-verification field calls the curve explicit, but the supplied text names it only as \(E\); this missing identity is an unresolved formulation issue.
\subsection{Short English statement}
For the intended elliptic curve, classify cubic radical fields that add points, or, if intended, increase its Mordell--Weil rank. The curve itself is not identified, and adding points is not automatically the same as increasing rank.
\subsection{Sources}
\begin{itemize}
\item Local UnsolvedMath registry, permanent ID 30000276, catalogue entry, \emph{Rank Growth over Cubic Radical Extensions}.
\item \emph{Explicit Methods in Number Theory} (2006), Oberwolfach Reports, pp. 1799--1866, DOI \url{https://doi.org/10.4171/owr/2005/32}.
\item Supplied reference labeled ``Dokchitser noncommutative Iwasawa computations'': \url{https://arxiv.org/abs/math/0509286}.
\item Supplied literature triage checked 2026-08-21; no defining equation for \(E\), rank theorem, or classification criterion is reproduced.
\end{itemize}
\end{document}
